\input format.tex

%\usepackage{arydshln}

\begin{document}

%%%
\setlength{\arrayrulewidth}{0.5pt}
\color{gray2}
\vspace*{0mm}
\fontsize{9.3pt}{9.8pt}\selectfont
\begin{tctabularx}{tabularx={p{15.4cm}<{\centering}p{0pt}}}
\\[-6.769pt]
\cellcolor{kelly} \raisebox{1pt}{\color{black70}\xiaosihao\lantxh\bf 肠道调节方案} \\
\end{tctabularx}

{
\lantxh
\vspace*{-4.25mm}
\noindent
\begin{longtable}{|@{~~~}p{15.35cm}|}
\hline
\parbox[c]{\hsize}{\vskip7pt\xiaowuhao {\lantxh\bf 补充益生元、益生菌}\\1.每天服用适量低聚半乳糖或低聚木糖产品，促进体内双歧杆菌、乳杆菌等益生菌的生长，提高肠道菌群多样性。

2.每天服用适量益生菌产品（如酸奶、含活菌的益生菌饮料等），增加肠道内益生菌的数量，有利于肠道菌群平衡。

注：如患炎症性肠病、慢性腹泻，不宜服用菊粉及低聚果糖，建议选择低聚半乳糖。
\vskip7pt}\\
\hline
\end{longtable}
}


\fontsize{9.3pt}{9.8pt}\selectfont
\begin{tctabularx}{tabularx={p{15.4cm}<{\centering}p{0pt}}}
\\[-6.769pt]
\cellcolor{kelly} \raisebox{1pt}{\color{black70}\xiaosihao\lantxh\bf 运动方案} \\
\end{tctabularx}

{
\lantxh
\vspace*{-4.25mm}
\noindent
\begin{longtable}{|p{1.25cm}@{}p{2cm}|p{11.732cm}|}
\hline
\parbox[c]{\hsize}{\vskip7pt\raisebox{-1.5ex}{\includegraphics[width=.9cm]{yundong.pdf}}\vskip7pt} & \parbox[c]{\hsize}{\centering\xiaowuhao\vskip7pt{运动}\vskip7pt} &
\parbox[c]{\hsize}{\vskip7pt\xiaowuhao {1.养成每天运动的好习惯，可选择竞技性运动或运动量较大的健身运动，如各种球类、跑步、爬山、骑自行车、游泳等。学习或工作较忙者，可利用零散时间进行以下运动，如快步走、爬楼梯、哑铃、深蹲、俯卧撑、跳绳等。

2.建议每天运动40分钟以上，每周5-7天。尽量减少静坐时间，每间隔1小时起来活动2\textasciitilde 3分钟。}\vskip7pt}\\
\hline
\end{longtable}
}
\end{document}
